% ============================================================
% discussion.tex
% ============================================================

This chapter synthesizes our empirical results, connects them
to the broader literature and policy landscape, and
establishes the identification limits that constrain
interpretation. We organize the discussion around five
substantive themes: the
PSLRA as a time-varying procedural filter, the decomposition
of the raw PSLRA association via composition adjustment, judicial
geography as a structural determinant of outcomes, the role
of case-level characteristics, and the identification limits
that constrain causal interpretation of all estimates.

\section{Summary of the Research Question and Approach}

This thesis set out to answer a deceptively simple question:
what determines whether a federal securities class action
settles, is dismissed, and how quickly? Answering it
correctly requires confronting two statistical challenges
that prior work has not simultaneously addressed. First,
settlement and dismissal are competing risks: they are
mutually exclusive terminal events, and studying settlement
duration conditional on settlement occurring recovers a
distorted picture of litigation dynamics. Second, the
dominant covariate of interest---PSLRA regime---has a
time-varying effect whose procedural mechanisms operate
through the motion-to-dismiss phase, implying a hazard ratio
that changes over litigation duration.

We address both challenges using a competing-risks survival
framework applied to 12,968 federal securities class actions
from the FJC Integrated Database. The framework combines
Aalen-Johansen nonparametric CIF estimation, cause-specific
Cox models, Fine-Gray subdistribution models,
inverse-probability-of-treatment-weighted (IPTW) models for
composition adjustment, and shared frailty models for
sensitivity to unobserved circuit-level heterogeneity. To
our knowledge, this is the first application of a full
competing-risks framework---incorporating both
composition-adjusted and random-effects extensions---to
population-level administrative data on federal securities
litigation.

\section{Policy Implications of the PSLRA Association}

Our results provide the most detailed empirical
characterization to date of the association between the
Private Securities Litigation Reform Act of 1995 and
securities litigation dynamics. The three-phase temporal decomposition
(Section~\ref{sec:piecewise}) is the single most
policy-relevant finding in this thesis. In the
first year, post-PSLRA cases are associated with a 95\% higher
instantaneous dismissal hazard (HR $= 1.948$, 95\% CI:
1.675--2.266). This is the PSLRA's heightened pleading
standard and automatic discovery stay operating as designed:
the ``strong inference'' of scienter requirement raises the
bar for surviving a motion to dismiss, and the discovery stay
prevents plaintiffs from using early discovery to supplement
insufficiently pleaded complaints.

The null association in years one to two (HR $= 1.050$,
$p = 0.582$) is equally important. Cases that survive the
first-year gauntlet face no differential dismissal risk
compared to pre-PSLRA survivors in the following year. This
pattern is consistent with the PSLRA acting as a one-time
filter at the pleading stage: once a case has demonstrated
sufficient pleading adequacy to survive initial review, the
reform no longer systematically disadvantages plaintiffs.
The re-emergence of a modest PSLRA association beyond two
years (HR $= 1.329$, $p < 0.001$) may reflect differences
in summary judgment practice or a selection effect among
long-duration post-PSLRA survivors.

From a policy standpoint, the three-phase temporal pattern
is consistent with the interpretation that the PSLRA
functions primarily as an early-stage procedural filter. The
first-year dismissal association is large and precisely
estimated; the subsequent null period suggests that surviving
cases proceed on approximately equal footing with pre-PSLRA
cases. However, we cannot determine from administrative data
alone whether the filtered cases are genuinely meritless or
whether some meritorious claims are caught in the same net.
Answering that question requires merging our dataset with
settlement amount data and case-specific merit proxies, and
remains an open empirical question that the identification
limits documented in our robustness analysis (particularly
the placebo test failure) make especially pressing.

The robustness analysis (Table~\ref{tab:robustness}) provides
important bounds on the settlement association, but also
reveals meaningful fragility.
Under the most liberal coding scheme (Scheme~B), the PSLRA
settlement HR rises from 0.445 to 0.722, still representing
an approximately 28\% reduction in settlement hazard. The
range across coding schemes ($[0.445, 0.722]$) depends on the
fraction of Code~12 voluntary dismissals that represent
private settlements. Under flexible time-trend controls
(a natural spline with $df{=}3$), the settlement HR
attenuates to 0.610 ($p = 0.001$)---still significant, but
substantially reduced from the raw estimate. Critically,
the settlement association loses statistical significance
entirely within individual circuits (Second Circuit:
HR $= 0.980$, $p = 0.953$; Ninth Circuit: HR $= 0.869$,
$p = 0.129$), and the 1992 placebo test detects significant
pre-existing trends before the reform was enacted (HR $= 0.719$,
$p < 0.001$), further limiting attribution to the PSLRA
specifically. The settlement finding is therefore directionally
consistent across all specifications but fragile to
circuit-specific decomposition and confounded by pre-existing
temporal trends.

\section{Composition-Adjusted Decomposition}

The IPTW triangulation (Section~\ref{sec:results_iptw})
addresses a natural question: how much of the raw PSLRA
association persists after adjusting for observable case-mix
changes, and how much is explained by those compositional
shifts? The answer
differs sharply by outcome. For dismissal, the
composition-adjusted MSM estimate (HR $= 1.526$) is
\emph{slightly amplified} relative to the unadjusted
HR of 1.466, indicating that concurrent compositional
shifts---changes in circuit filing patterns, statutory basis
mix, and MDL prevalence---actually worked \emph{against} the
dismissal association. After adjusting for these observable
changes in case composition, the PSLRA's association with
dismissal is modestly larger, not smaller. The full
triangulation (unadjusted 1.466, regression-adjusted 1.746,
doubly robust 1.804, MSM 1.526) confirms that the dismissal
finding is robust to the weighting strategy chosen.

For settlement, composition adjustment attenuates the raw
association substantially: from a baseline HR of 0.439 to
an MSM HR of 0.690 (and regression-adjusted 0.701).
Approximately half the raw settlement suppression is
explained by observable shifts in case composition across
the pre- and post-PSLRA periods. This does not mean the
remaining association is causally attributable to the reform:
the IPTW framework adjusts for observed compositional
differences but cannot control for unobserved secular
trends in judicial behavior, plaintiff strategy, or economic
conditions.

The composition-adjusted estimates inherit two important
caveats. First, several covariates included in the
propensity score model---circuit, origin, MDL status, and
statutory basis---may themselves be consequences of the
PSLRA rather than pre-treatment confounders. If the reform
altered where and how cases were filed, the IPTW adjustment
absorbs some of the reform's \emph{indirect} effects,
producing a lower bound on the total PSLRA association.
Second, the 1992 placebo test (dismissal HR $= 0.719$,
$p < 0.001$) detects significant pre-existing trends in the
pre-PSLRA period, indicating that the model cannot cleanly
separate the reform from concurrent temporal drift. These
estimates are best understood as a decomposition---isolating
the component of the raw association not explained by
observable case-mix changes---rather than a causal
identification.

\section{Circuit Heterogeneity and Judicial Geography}

The dominance of circuit geography in our results---confirmed
by both the cause-specific Cox models and the Fine-Gray
subdistribution analysis---raises fundamental questions
about the role of judicial institutions in shaping litigation
outcomes. The Second Circuit's anomalous position is the
most striking geographic finding. Every other circuit has a
settlement hazard 4.1 to 17.6 times that of the Second
Circuit (Table~\ref{tab:cox_ext_settle}), yet the Second
Circuit has the highest cumulative
dismissal incidence. This combination defines a litigation
environment distinctively hostile to monetary recovery for
plaintiffs' counsel. Importantly, our data reflect the full
Second Circuit (including EDNY, NDNY, WDNY, and the
District of Vermont), so this pattern is not attributable
solely to the Southern District of New York, though SDNY
contributes the largest share of cases.

Several mechanisms could explain this pattern. The Second
Circuit handles a disproportionate share of the most complex
and high-stakes securities cases, and defendants in these
cases are typically represented by elite defense counsel
with substantial expertise, which may reduce early settlement
probability. Second Circuit judges also have extensive
experience with securities class actions and may apply
demanding standards at class certification, reducing the
expected value of cases to plaintiffs. The PSLRA $\times$
Circuit interaction analysis adds a temporal dimension:
within the Second Circuit (the reference category), the PSLRA
is associated with a 142\% higher dismissal hazard
(HR $= 2.422$, $p < 0.001$), while the reform's settlement
association is non-significant (HR $= 1.344$, $p = 0.371$),
consistent with the Second Circuit's already-restrictive
settlement environment leaving less room for additional
suppression.

The most striking interaction attenuations for dismissal are
in the Fourth Circuit (HR $= 0.266$, $p < 0.001$) and
Eleventh Circuit (HR $= 0.441$, $p < 0.001$), indicating
that the PSLRA's dismissal amplification is strongly
attenuated in these jurisdictions. Prior work by
\citet{johnson2001} and \citet{choi2007} has documented
circuit-level variation in PSLRA application; our
competing-risks framework provides the first quantitative
estimates of how this variation manifests in the joint
distribution of settlement and dismissal timing.

The shared frailty models (Section~\ref{sec:frailty}) add an
important dimension to this geographic analysis. The settlement frailty variance is
substantial ($\hat{\theta} = 0.449$), indicating that
unobserved circuit-level factors---judicial culture, local
bar norms, case-management practices not captured by the
observed covariates---contribute meaningfully to settlement
heterogeneity beyond what the fixed-effect circuit indicators
capture. By contrast, the dismissal frailty variance is
modest ($\hat{\theta} = 0.033$), suggesting the fixed-effect
circuit indicators already capture most circuit-level
variation in dismissal timing. Importantly, the PSLRA
hazard ratio is virtually identical between fixed-effect
and random-effect models (difference $< 0.001$), which is
mathematically expected for a national-level treatment
applied uniformly across all circuits. The frailty model's
contribution is not an alternative PSLRA estimate but rather
a quantification of the residual geographic heterogeneity
and cluster-robust standard errors that widen confidence
intervals by a factor of 2.8$\times$ for settlement
($p = 0.023$) and 2.0$\times$ for dismissal ($p < 0.001$).

The policy implication is sobering. If the PSLRA's
filtering effects vary substantially across circuits, then
plaintiffs filing in different jurisdictions face
qualitatively different litigation environments. This
raises questions about regulatory forum shopping, the
appropriate geographic scope of securities law reform, and
whether the PSLRA inadvertently created
jurisdiction-specific barriers to recovery.

\section{Structural Case Characteristics: MDL, Origin,
and Statutory Basis}

The MDL finding (Section~\ref{sec:results_characteristics})
is methodologically novel. MDL consolidation is associated
with large bilateral suppression
of both cause-specific hazards---approximately 72\% for
settlement (HR $= 0.284$, 95\% CI: 0.212--0.379) and 75\%
for dismissal (HR $= 0.247$, 95\% CI:
0.219--0.277)---indicating that MDL cases resolve at
dramatically lower instantaneous rates regardless of their
eventual outcome. The natural interpretation is that
consolidation introduces coordination costs and procedural
delays---bellwether trials, coordinated discovery schedules,
global resolution frameworks---that slow individual case
resolution. The Fine-Gray subdistribution analysis
complicates this picture in an important way: despite the
dramatically lower instantaneous hazards, MDL cases show
a modestly higher cumulative settlement probability
(SHR $= 1.167$, 95\% CI: 0.913--1.492, $p = 0.217$).
Because bilateral suppression of both hazards keeps large
numbers of MDL cases in the risk set throughout the study
period, cumulative settlement probability accumulates to a
similar level as non-MDL cases despite the substantially
lower instantaneous settlement rate---though this difference
is not statistically significant. This sign reversal is
visible only through competing-risks analysis and illustrates
why both cause-specific and subdistribution frameworks are
necessary for a complete picture of resolution dynamics.

Case origin results suggest that removed cases are associated
with lower dismissal hazard (HR $= 0.791$, 95\% CI:
0.632--0.990) and higher settlement hazard (HR $= 1.540$,
95\% CI: 1.053--2.252), consistent with the interpretation
that cases initially filed in state court involve claims
where settlement is a more natural resolution than federal
dismissal. Statutory basis results are subtler: Section~11
cases are associated with a modestly higher dismissal hazard
than Section~10(b) cases (HR $= 1.118$, 95\% CI:
1.044--1.197) and a modestly higher settlement hazard
(HR $= 1.167$, 95\% CI: 1.013--1.345), which may reflect
tracing and standing requirements in IPO and secondary
offering contexts that create both additional dismissal
opportunities and, for surviving cases, stronger settlement
incentives.

\section{Limitations}

Several limitations constrain interpretation of our
findings. We organize them by severity, beginning with
those that most directly affect the credibility of our
central estimates.

\paragraph{Identification limits.}
The 1992 placebo test is the single most important
limitation. By placing a fictitious reform at 1992 within
the pre-PSLRA period (1990--1995), we detect a significant
dismissal HR of 0.719 ($p < 0.001$), indicating that the
model captures pre-existing trends in dismissal rates
\emph{in the absence of any legislative change}. This result
precludes clean causal attribution of the observed PSLRA
association to the reform in isolation and strengthens our
decision to frame all estimates as associational or
composition-adjusted rather than causally identified. A
narrow-window analysis (1993--1998) confirms that the
dismissal association survives in a restricted band
(HR $= 1.394$, $p < 0.001$), but this is not a formal
regression discontinuity design and cannot fully address the
placebo concern.

\paragraph{Post-treatment bias in composition adjustment.}
Several covariates included in the IPTW propensity score
model---circuit filing location, case origin, MDL
consolidation, and statutory basis---may themselves be
consequences of the PSLRA rather than pre-treatment
confounders. If the reform altered where and how securities
class actions were filed, the composition-adjusted estimates
absorb some of the reform's \emph{indirect} effects,
producing a lower bound on the total PSLRA association.
The composition-adjusted HR should therefore be interpreted
as the component of the raw association not explained by
observable case-mix shifts, not as the ``true'' reform
effect.

\paragraph{Unobserved heterogeneity.}
Our hazard ratios represent population-averaged associations
conditional only on observed administrative covariates.
Prior literature identifies case-specific merit
signals---such as the magnitude of stock price drops, the
presence of SEC investigations, or accounting
restatements---as strong predictors of outcomes
\citep{cox2008, johnson2001}. Because these metrics are
absent from our dataset, our models cannot isolate how much
of the PSLRA's early dismissal association is driven by the
dismissal of genuinely meritless claims versus structurally
complex but valid claims. Litigation outcomes are also
highly dependent on the specific judge assigned to a case
and the caliber of the law firms involved
\citep{wang2022}. While we implement circuit-level shared
frailty models, the IDB lacks consistent judge identifiers
or attorney data; finer-grained clustering at the judge or
counsel level remains unmodeled.

\paragraph{Small-cluster frailty caveat.}
The shared frailty models operate with only 11 circuit-level
clusters, well below the 20--30 clusters typically
recommended for reliable variance estimation. The frailty
variance estimates ($\hat{\theta}_{\text{settle}} = 0.449$,
$\hat{\theta}_{\text{dismiss}} = 0.033$) should be treated
as indicative rather than precise, and may be biased
downward.

\paragraph{Disposition coding and Code~6 disaggregation.}
Our primary disposition scheme disaggregates FJC Code~6
(``Judgment on Motion Before Trial'') using the IDB
JUDGMENT field, reclassifying plaintiff victories as
settlements and defendant victories as dismissals. Of the
701 plaintiff-victory reclassifications, 93.4\% are
post-PSLRA cases, introducing a potential era confound: if
post-PSLRA courts are more likely to record motion outcomes
via Code~6 than pre-PSLRA courts, the reclassification may
amplify the apparent PSLRA association. The robustness of
our results across all three coding schemes (A, B, C)
provides partial reassurance, but this coding asymmetry
should be noted when interpreting the precise magnitudes.

\paragraph{Proportional hazards violations.}
Grambsch-Therneau global tests reject proportional hazards
in both extended models (settlement: $\chi^2 = 301.1$,
$p < 0.001$; dismissal: $\chi^2 = 471.2$, $p < 0.001$).
All constant hazard ratio estimates in this thesis must
therefore be interpreted as time-averaged summaries of a
more complex temporal relationship, not as constant
instantaneous risks. The piecewise Cox model partially
addresses this for the PSLRA covariate, but other
covariates' time-varying effects remain unmodeled.

The limitations discussed above---particularly the
identification constraints from the placebo test and the
unobserved heterogeneity that our administrative data cannot
capture---motivate the research extensions proposed in
Chapter~\ref{ch:conclusion}, which also synthesizes the thesis's
central contributions and their implications for empirical
legal scholarship.